\chapter{Game Theory}

\section{Strategic-Form Games}
\begin{itemize}
    \item Game: players, available strategies, and payoffs for each strategy profile
    \item Strategy profile: one strategy chosen by each player
    \item Payoff: value or outcome a player receives from a strategy profile
    \item Best response: strategy that maximizes a player's payoff given the other players' strategies
    \item Strict best response: strategy that gives a strictly higher payoff than every alternative, given the other players' strategies
    \item More than one best response possible when strategies tie for the highest payoff
\end{itemize}

\section{Dominant Strategies}
\begin{itemize}
    \item Dominant strategy: best response to every possible strategy profile of the other players
    \item Strictly dominant strategy: gives a strictly higher payoff than every alternative for every possible strategy profile of the other players
    \item Strictly dominant strategy: rational choice regardless of what the other players do
    \item A game may have no dominant strategy, or may have dominant strategies for some or all players
\end{itemize}

\section{Example Games}
\subsection{Presentation vs. Exam}
\begin{itemize}
    \item Compare each player's available choices and payoffs in every outcome
    \item Mark each player's best response to each choice of the other player
    \item Use the payoff matrix to determine whether choices align or conflict
\end{itemize}

\subsection{Prisoner's Dilemma}
\begin{itemize}
    \item Two players independently choose to cooperate or defect
    \item Defecting: strictly dominant strategy for each player
    \item (Defect, Defect): Nash equilibrium, despite mutual cooperation giving both players a better outcome
    \item Individual incentives can produce an outcome worse for both players than cooperation
\end{itemize}

\subsection{A New Product}
\begin{itemize}
    \item Analyze the available choices and payoffs for each player using the lecture payoff matrix
    \item Check how one player's best response changes with the other player's choice
    \item Use best responses to identify any Nash equilibria
\end{itemize}

\subsection{Three Clients}
\begin{itemize}
    \item Identify the players, available strategies, and payoffs from the lecture example
    \item For each player, compare payoffs across their choices while holding the other players' choices fixed
    \item Mark mutual best responses to locate Nash equilibria
\end{itemize}

\section{Nash Equilibrium}
\begin{itemize}
    \item Two-player strategy profile $(S,T)$: Nash equilibrium if $S$ is a best response to $T$ and $T$ is a best response to $S$
    \item General case: every player's strategy is a best response to the strategies chosen by all other players
    \item At equilibrium: no player can improve their payoff by changing strategy alone
    \item Finite-game existence theorem: every game with finitely many players and finitely many strategies has at least one Nash equilibrium in mixed strategies
    \item Pure-strategy Nash equilibrium: equilibrium using a single strategy per player; not guaranteed to exist
\end{itemize}

\section{Focal Points}
\begin{itemize}
    \item Focal point: salient outcome that helps players coordinate without communication
    \item Salience may come from conventions, shared expectations, or a distinctive feature of the game
    \item Focal points can help players select among multiple coordination-game equilibria
\end{itemize}

\section{Mixed Strategies}
\begin{itemize}
    \item Mixed strategy: probability distribution over a player's available pure strategies
    \item Pure strategy: one action chosen with probability $1$
    \item Expected payoff: probability-weighted average payoff across possible outcomes
    \item For strategy profile $(s_i,s_{-i})$ drawn from mixed strategies:
          \[
              \mathbb{E}[u_i] = \sum_{s_i,s_{-i}} \Pr(s_i)\Pr(s_{-i})u_i(s_i,s_{-i})
          \]
    \item Mixed-strategy Nash equilibrium: each player's probability distribution is a best response to the others' distributions
    \item In equilibrium, every pure strategy used with positive probability gives the player the same expected payoff
\end{itemize}

\subsection{Finding a Mixed-Strategy Equilibrium}
\begin{itemize}
    \item Assign probabilities to one player's strategies
    \item Choose probabilities that make the other player indifferent among their pure strategies
    \item Repeat in the opposite direction to find the first player's equilibrium probabilities
    \item Indifference makes mixing optimal; a player would otherwise put all probability on a higher-payoff pure strategy
\end{itemize}

\section{Matching Pennies}
\begin{itemize}
    \item Each player chooses Heads or Tails
    \item Player 1 wins when choices match; Player 2 wins when choices differ
    \item No pure-strategy Nash equilibrium: either player can switch and improve their payoff
    \item Let $q$ be the probability Player 2 chooses Heads
    \item Player 1's expected payoff from Heads: $q - (1-q) = 2q-1$
    \item Player 1's expected payoff from Tails: $(1-q)-q = 1-2q$
    \item Player 1 is indifferent when $2q-1=1-2q$, so $q=\frac{1}{2}$
    \item By symmetry, Player 1 also chooses Heads with probability $\frac{1}{2}$
    \item Mixed-strategy Nash equilibrium: each player chooses Heads and Tails with probability $\frac{1}{2}$
\end{itemize}

\section{Coordination Game}
\begin{itemize}
    \item Players benefit from choosing strategies that match or coordinate
    \item Coordination games often have multiple Nash equilibria
    \item Equilibrium selection: players need expectations, conventions, or focal points to coordinate on one equilibrium
    \item A coordinated outcome may be better for everyone than a mismatched outcome
\end{itemize}

\section{Network Traffic}

\begin{itemize}
    \item Players choose routes in a network to minimize their own travel time
    \item Travel time on a route depends on the number of players using it
    \item Nash equilibrium occurs when no player can reduce their travel time by unilaterally changing routes
    \item Equilibria may be inefficient compared to a centrally coordinated solution (Braess's paradox)
\end{itemize}